%! Linear Transformations and Matrices — Unit 4, Lesson 4.12 — Group Activity (Answer Key)
\documentclass[10pt]{article}
\usepackage{precalculus-article}
\usepackage{precalculus-key}   % pulls in -boxes; do NOT also load -boxes
\newcommand{\TallMath}[1]{$\displaystyle #1\rule[-1.4em]{0pt}{3.2em}$}

\begin{document}

\pageheader{Unit 4, Lesson 4.12}{Group Activity --- Answer Key}
\namepartnerperiod

\vspace{0.1in}

\begin{scenariobox}[Exploring Geometric Transformations --- Key]{sky}
\small
Each tier below uses a different $2\times 2$ matrix.  Your goal is to compute output
vectors and connect the matrix to a geometric action.
\end{scenariobox}

\vspace{0.1in}

%----------------------------------------------------------------------
% Tier R Key
%----------------------------------------------------------------------
\begin{tcolorbox}[colback=white, colframe=black!40,
    title=\textbf{Tier R --- Reinforce --- Key},
    fonttitle=\bfseries, arc=2mm, left=3mm, right=3mm, top=2mm, bottom=2mm]
\small
Let $A = \begin{bmatrix}2 & 0 \\ 0 & 2\end{bmatrix}$.

\textbf{(a)}
$A\vec{v}_1 = \begin{bmatrix}2&0\\0&2\end{bmatrix}\begin{bmatrix}3\\1\end{bmatrix}
= $ \ans{$\begin{bmatrix}6\\2\end{bmatrix}$} \qquad
$A\vec{v}_2 = \begin{bmatrix}2&0\\0&2\end{bmatrix}\begin{bmatrix}-1\\4\end{bmatrix}
= $ \ans{$\begin{bmatrix}-2\\8\end{bmatrix}$}

\vspace{0.10in}
\textbf{(b)} $A\begin{bmatrix}0\\0\end{bmatrix} = $ \ans{$\begin{bmatrix}0\\0\end{bmatrix}$}

\vspace{0.08in}
\textbf{(c)} \ans{This transformation scales every input vector by a factor of 2 (dilation from the origin with scale factor 2).}
\end{tcolorbox}

\vspace{0.1in}

%----------------------------------------------------------------------
% Tier A Key
%----------------------------------------------------------------------
\begin{tcolorbox}[colback=white, colframe=black!40,
    title=\textbf{Tier A --- At Standard --- Key},
    fonttitle=\bfseries, arc=2mm, left=3mm, right=3mm, top=2mm, bottom=2mm]
\small
Let $A = \begin{bmatrix}0 & -1 \\ 1 & 0\end{bmatrix}$.

\textbf{(a)} Input matrix $= $ \ans{$\begin{bmatrix}1&0&-1\\0&1&-1\end{bmatrix}$}

\vspace{0.10in}
\textbf{(b)} Output matrix:

$\begin{bmatrix}0&-1\\1&0\end{bmatrix}\begin{bmatrix}1&0&-1\\0&1&-1\end{bmatrix}
= $ \ans{$\begin{bmatrix}0&-1&1\\1&0&-1\end{bmatrix}$}

\vspace{0.10in}
\textbf{(c)} Image vertices: \ans{$(0,1)$, $(-1,0)$, $(1,-1)$}

\vspace{0.08in}
\textbf{(d)} \ans{This transformation rotates every vector $90°$ counterclockwise about the origin.}
\end{tcolorbox}

\vspace{0.1in}

%----------------------------------------------------------------------
% Tier E Key
%----------------------------------------------------------------------
\begin{tcolorbox}[colback=white, colframe=black!40,
    title=\textbf{Tier E --- Extend --- Key},
    fonttitle=\bfseries, arc=2mm, left=3mm, right=3mm, top=2mm, bottom=2mm]
\small
Let $A = \begin{bmatrix}1&2\\0&1\end{bmatrix}$,
$\vec{v}=\begin{bmatrix}1\\2\end{bmatrix}$, $\vec{w}=\begin{bmatrix}3\\-1\end{bmatrix}$.

\textbf{Step 1.}
$\vec{v}+\vec{w} = \begin{bmatrix}4\\1\end{bmatrix}$. \quad
$L(\vec{v}+\vec{w}) = A\begin{bmatrix}4\\1\end{bmatrix} = $ \ans{$\begin{bmatrix}6\\1\end{bmatrix}$}

\vspace{0.08in}
\textbf{Step 2.}
$L(\vec{v}) = A\begin{bmatrix}1\\2\end{bmatrix} = $ \ans{$\begin{bmatrix}5\\2\end{bmatrix}$};
\quad $L(\vec{w}) = A\begin{bmatrix}3\\-1\end{bmatrix} = $ \ans{$\begin{bmatrix}1\\-1\end{bmatrix}$}

$L(\vec{v})+L(\vec{w}) = \begin{bmatrix}5\\2\end{bmatrix}+\begin{bmatrix}1\\-1\end{bmatrix}
= $ \ans{$\begin{bmatrix}6\\1\end{bmatrix}$}

\vspace{0.08in}
\textbf{Step 3.} \ans{The results are equal. This confirms that $L(\vec{v})=A\vec{v}$ satisfies the additive property, which is required for a linear transformation.}
\end{tcolorbox}

\begin{teachernote}
Tier A: the rotation matrix $\begin{bmatrix}0&-1\\1&0\end{bmatrix}$ is the standard
$90°$ CCW rotation matrix --- worth naming explicitly in debrief.
Tier E: students may struggle setting up the computation; prompt them to compute
$A(\vec{v}+\vec{w})$ first, then compare.
\end{teachernote}

\end{document}
